\documentclass[11pt,a4paper]{article}
\usepackage[T1]{fontenc}
\usepackage[utf8]{inputenc}
\usepackage{amsmath,amssymb}
\usepackage{geometry}
\geometry{margin=1.1in}
\usepackage{graphicx}
\graphicspath{{./}{../dossier/figures/}}
\usepackage{booktabs}
\usepackage{float}
\usepackage[expansion=false]{microtype}
\usepackage[colorlinks=true,linkcolor=black,urlcolor=blue!50!black,citecolor=black]{hyperref}
\usepackage{xcolor}

\title{Same Choices, Different Reasonings:\\ A Reservation-Price Design for Selection Neglect and Elicitation Effects}
\author{Francesco Pinna \and Emanuele Sebastianelli}
\date{November 2026\thanks{Course project for Experimental and Behavioural Economics, LUISS Guido Carli, taught by Daniela Di Cagno and Irene Maria Buso, academic year 2026--2027. The design is anchored to the public version of the companion preprint, Pinna (2026), and theorem, proposition and equation numbers cited without further qualification refer to that version.}}

\begin{document}
\maketitle

\begin{abstract}
\noindent A buyer who offers too much for a good of uncertain quality may value the good too highly or may neglect the selection operated by the seller, and the final offer alone does not separate the two explanations. We present an experimental design in which each participant states reservation prices for the same good against two programmed sellers, one always available and one more often available when the good is worth little, so that the gap between the two prices measures the perceived selection while the first price measures the perceived value. Three elicitation protocols that differ only in when prices are collected turn measurement itself into a treatment. The design implements the process measurements requested by the companion preprint on dynamic identification of behavioral reasoning, whose theory explains both why terminal choices cannot separate reasoning channels and how to measure without distorting what is measured. Two preregistered hypotheses ask whether the average degree of selection neglect is positive and whether final prices are equivalent, within three points, between the lightest and the richest protocol, and a third question, examined only where neglect is found, asks which of two descriptions of it fits the prices better. Every number reported here is a design calculation or a simulation of the analysis procedure, reproducible from the programs that travel with the design package.
\end{abstract}

\section{Introduction and research question}

A final action contains information about the rule that generated it, while it usually reveals much less about the path along which that rule was implemented. In a purchase decision this ambiguity takes a concrete form, since an offer above the rational benchmark can come from an optimistic valuation of the good or from neglect of the fact that the seller's willingness to trade is itself informative. The two mechanisms can produce the same offer, and they make different predictions for intermediate beliefs and for the response to new information, which is the sense of the title of this project.

The design asks whether the gap between the reservation price for an always-available seller and the reservation price for a selective seller is smaller than the gap implied by correct conditioning, once the individual valuation of the good has been taken into account, and whether collecting intermediate decisions changes the final price, so that the elicitation technology is studied as a treatment whose effect is measured. Where the neglect is found, a further question concerns its form, since a buyer who averages the probability she held before the sale with the correct one and a buyer who discounts the evidence carried by the sale state nearly the same prices under moderate selection and different prices under extreme selection. The dynamic bridge of the preprint stays outside the confirmatory part.

The questions on the neglect and on the effect of measurement come from the companion preprint, Pinna (2026), which develops a joint theory of endpoint and process identification in finite games and asks, in its Section 10.3, in its Section 12.4 and in its conclusions, for an experiment with randomized report timing, with a passive arm and with repeated reports against an unchanged information set from which response noise can be estimated. The present design is that experiment in purchase form, as Section 3 makes precise. The protocol dossier that accompanies this document contains the full specification, the participant instructions, the analysis plan and the verification programs, and the explanatory note in the same folder presents the design to a reader of experimental economics with a numerical example that runs through the text.

\section{Entry points for the reader}

We make the entry points into the problem and into the method explicit, so that a reader can reach the design through a short and familiar route rather than through the full preprint.

For the problem, the selection structure of the task is the market for lemons of Akerlof (1970), and its classical laboratory form is the acquiring-a-company problem of Samuelson and Bazerman (1985), whose field counterpart is the winner's curse documented by Kagel and Levin (1986). Charness and Levin (2009) strip the problem of strategic interaction and show that the overbidding survives against a programmed counterpart, which is the simplification adopted here, while Esponda and Vespa (2014) trace the failure to the extraction of information from hypothetical events and Enke (2020) documents neglect of the selection behind the data one happens to see. Eyster and Rabin (2005) provide the equilibrium formalization through the cursed contraction of the perceived dependence between information and behavior, and the parameter of that contraction is the object our prices target. The closest experimental antecedent in the environment of this course is the acquiring-a-company experiment with value messages and information leaks of Di Cagno, G\"uth, Lohse, Marazzi and Spadoni (2024), from which the present design differs by replacing the human seller with a declared programmed rule whose sale does not respond to the price, so that the strategic forecast is removed, the benchmark becomes computable and the inference from the sale is isolated from the further step by which a buyer's own offer decides which goods she obtains.

For the method, Schlag, Tremewan and van der Weele (2015) survey belief elicitation and give the map on which this design takes a position, while Danz, Vesterlund and Wilson (2022) document systematic distortions of a binarized scoring rule that worsen when incentives are explained better, which is the reason the design asks for reservation prices and never for probabilities. The incentive logic of a price list with a randomly drawn row goes back to Becker, DeGroot and Marschak (1964), whose mechanism punishes small departures from the true reservation price very lightly, the flat maximum that Harrison (1989) pointed out for laboratory auctions, and is frequently mistaken for one in which the stated price is the price paid, as Cason and Plott (2014) document, two weaknesses to which Section 4.5 responds. Brandts and Charness (2011) survey the evidence that the elicitation format can itself change behavior, and our second hypothesis measures such an effect with the equivalence procedure of Schuirmann (1987). Enke and Graeber (2023) motivate the alternative of an attraction of the answers toward the center of the list, against which the design defends itself through lists with an announced value and a randomly shifted price list, and Grether (1980) is the reference for the description of the neglect as a discount on evidence.

\section{The theoretical core of the preprint}

\subsection{Two states and the cursed target}

The preprint separates two simplex-valued objects with distinct probabilistic roles. The rational posterior is the belief generated by the experimental information and remains a martingale, while the perceived belief is a second process whose drift records systematic distortion and whose quadratic variation records latent exploration. Distortion enters through the cursed transformation of Section 4.2 of the preprint, which acts on a joint probability table over a payoff state and another player's action, here the value of the good and the availability of the seller, and which mixes that table with the product of its marginals. Writing the joint table as $p$ and its two marginals as $p_{\Theta}$ and $p_{A}$, the transformation with contraction parameter $\chi$, a scalar in the unit interval, is
\[
\Gamma_{\chi}(p) \;=\; (1-\chi)\,p \;+\; \chi\, p_{\Theta}\otimes p_{A},
\]
so that marginals are preserved and only the association is contracted. Proposition 4.2 of the preprint shows that the transformation contracts mutual information at least linearly, and Proposition 4.3 shows that the induced drift of the perceived belief has a derivative in $\chi$ equal to the rational association gap, which is the pathwise content behind the price identity of Section 4.2 below.

\subsection{What terminal choices cannot reveal}

Theorem 9.1 of the preprint states that the score of any observed coarsening of the reasoning path is the conditional expectation of the entire-path score, so that path information decomposes exactly into observable information and a missing term. A direction that changes only the conditional path law given the observation carries zero observable information even when its path information is positive, which is the formal statement of endpoint impossibility and the reason additional process measurements are proposed at all. The preprint also defines, in its equation (77) and in its Section 13.1, the efficient information for a target parameter after profiling nuisance parameters, and its reanalysis of the Costa-Gomes and Crawford data shows that games with more efficient depth information sort subject types better in choices while generic lookup counts do not reproduce the ordering, so that process measurements help only when they are tied to the relevant score.

\subsection{Elicitation as treatment and as measurement}

Section 10.3 of the preprint writes a reported coordinate at elicitation time $\ell$ under schedule $D$ as its equation (94), in which the latent path may itself depend on the schedule, a report-specific displacement is added, and a transitory noise term closes the model. Repeated reports against an unchanged information set identify the noise variance, whereas a systematic elicitation effect cannot be absorbed into that variance, so that the preprint asks for randomized report timing and for a passive arm. The present design realizes that request term by term, since protocol F below is the passive arm, protocol R is the active arm with two intermediate measurements, protocol S is the intermediate treatment with a single measurement per trial, the lists with an announced value provide repeated price reports at unchanged information from which the response noise is estimated, and the second hypothesis is the test of the systematic elicitation effect of equation (94). The dossier implements the measurement theory of Sections 9 and 10 of the preprint, while the three-game battery of its Section 12 remains a distinct experimental venue, and the selective seller carries the state-action association studied in that battery into the form of a purchase decision with reservation prices.

\section{Design}

\subsection{The purchase task}

In every purchase trial the good is worth 30 or 90 points, the announced prior probability of the high value, a scalar denoted $\pi_{0}$, equals 0.4 or 0.6, and two certificates conditionally independent given the value arrive in random order, one correct with probability 0.60 and the other with probability 0.70, with symmetric accuracies that are always displayed. The participant faces two sellers for the same good. Their availability is generated by an announced programmed rule, it does not depend on prices or answers, and the probability of being available when the value is high and when it is low, two scalars in the unit interval denoted $a$ and $b$, takes the values in Table~\ref{tab:sellers}.

\begin{table}[H]
\centering
\begin{tabular}{lcc}
\toprule
Seller type & Available if value 30 & Available if value 90\\
\midrule
Always available & 1 & 1\\
Weak selection & 0.65 & 0.35\\
Strong selection & 0.80 & 0.20\\
Very strong selection & 0.95 & 0.05\\
Independent placebo & 0.50 & 0.50\\
\bottomrule
\end{tabular}
\caption{Availability rules of the programmed sellers.}
\label{tab:sellers}
\end{table}

The 22 purchase trials of a participant comprise twelve joint trials, which cross the two priors with the three intensities of selection in two repetitions and present both sellers on the same screen, six single-seller trials, two placebo trials and two revelation trials, in which the realized availability is shown after the second certificate and before the final list. In a single-seller trial the two members of a pair of participants see the same good, the same prior and the same certificates, and while one of them states a price only for the always-available seller the other states a price only for the selective seller, with roles that alternate so that each participant prices each seller alone three times. The members of a pair never interact and never see each other's answers, so that the pair is a unit of shared evidence and of randomization that involves no strategic relation. For each seller the participant indicates up to which price she would buy by typing a threshold in a field that shows the range of allowed prices, on a list with unit step whose format runs from 20 to 90 in half of the trials and from 30 to 100 in the other half, so that both formats have 71 rows, cover every economically possible price and differ only in their numerical center, which moves from 55 to 65 and allows the check on the attraction toward the center. The threshold fills the monotone list, the summary screen shows both regions before confirmation, nothing is preselected, and if a list is drawn for payment the computer draws one price and applies the stated policy at that price, so that the stated threshold is never automatically the price paid. Figure~\ref{fig:task} shows a facsimile of the screen.

\begin{figure}[H]
\centering
\includegraphics[width=0.9\linewidth]{facsimile_X.png}
\caption{Facsimile of the initial request. The boxes are empty and no threshold is preselected. The screen illustrates the task and is not a collection application.}
\label{fig:task}
\end{figure}

\subsection{Benchmarks and the price identity}

Write the current posterior probability of the high value, a scalar in the open unit interval, as $\pi$. Under the payment rule of Section 4.5, which converts points into the probability of a single prize, the reservation price of a buyer who conditions correctly is, for the always-available seller and for a selective seller with availabilities $a$ and $b$,
\[
U(\pi) \;=\; 30+60\pi,
\qquad
C_{a,b}(\pi) \;=\; 30+60\,\frac{a\pi}{a\pi+b(1-\pi)}.
\]
Under the cursed transformation of the preprint applied to the joint table of value and availability, the selective price is exactly
\[
S(\pi,\chi) \;=\; C_{a,b}(\pi) + \chi\,\bigl[U(\pi)-C_{a,b}(\pi)\bigr],
\]
a straight line between the two benchmarks, in which $\chi$ is read as the degree of selection neglect, equal to zero for a buyer who draws the correct inference from the sale and to one for a buyer who ignores it. If the first price measures the subjective posterior through $\widetilde{\pi}=(U-30)/60$ and the announced conditional probabilities are represented correctly, the degree recovered from a pair of prices is
\[
\widehat{\chi} \;=\; \frac{S-C_{a,b}(\widetilde{\pi})}{\,U-C_{a,b}(\widetilde{\pi})\,}.
\]
The coefficient so recovered is an effective quantity, which coincides with the parameter of the preprint in the static restriction with a correct posterior, while the dynamic model of the preprint does not force the perceived process to sit on its target at every instant. The identifying gap $U(\pi)-C_{a,b}(\pi)$ equals $60(b-a)\pi(1-\pi)/\{a\pi+b(1-\pi)\}$, it is positive for negative selection at interior posteriors, and the map from posterior and degree to the two prices has Jacobian determinant equal to sixty times that gap, which gives local rank two in the static restriction. At the posterior 0.4 the correct prices are 54.00 for the always-available seller and 45.85, 38.57 and 32.03 for the selective seller under weak, strong and very strong selection, and a degree of 0.25 moves the three selective prices to 47.89, 42.43 and 37.53.

\subsection{Joint statements and the form of the neglect}

A participant who states two prices on the same screen may keep them close for reasons unrelated to the selection, and if the price for the selective seller moves toward the other by a share $\kappa$ of the distance between them, a share called the pull, the joint trials alone identify only the composite $\chi+\kappa(1-\chi)$, in which a positive pull is read as neglect. The single-seller trials remove the confusion, because the pull cannot operate between prices stated by different persons, so that the gap between partners reflects the neglect alone while the gap within persons reflects neglect and pull together, and the likelihood of Section 5.1 uses both kinds of trial with the pull as a common parameter.

The price identity of Section 4.2 describes a buyer who averages the probability she held before the sale with the correct one, and a second description, in the tradition of Grether (1980), has the buyer discount the evidence of the sale on the scale of the log-odds. Writing the correct probability after a sale as $\pi^{\mathrm{sale}}=a\pi/\{a\pi+b(1-\pi)\}$, the probabilities that the two descriptions assign to the high value after a sale, denoted $\pi_{\mathrm{A}}$ and $\pi_{\mathrm{D}}$, satisfy
\[
\pi_{\mathrm{A}} \;=\; (1-\chi)\,\pi^{\mathrm{sale}}+\chi\,\pi,
\qquad
\log\frac{\pi_{\mathrm{D}}}{1-\pi_{\mathrm{D}}} \;=\; \log\frac{\pi}{1-\pi}+(1-\chi)\log\frac{a}{b}.
\]
The two rules, called averaging and dampening, coincide without neglect and with full neglect, stay close under weak selection and part ways under extreme selection, which is the reason for the third intensity, since at the posterior 0.4 and at a degree of 0.5 the selective price exceeds the correct one by 4.1, 7.7 and 11.0 points across the three intensities under the averaging rule and by 3.9, 6.4 and 5.9 points under the dampening rule. The analysis connects the two descriptions through a weight between zero and one, called the rule weight and common to all participants, which multiplies the difference between $\pi_{\mathrm{D}}$ and $\pi_{\mathrm{A}}$ added to $\pi_{\mathrm{A}}$. Both rules solve the problem of a buyer who looks for a single belief as close as possible, in relative entropy, to the correct one and to the one that ignores the sale, and they correspond to the two orders of the arguments of that distance, the averaging rule being the one that the preprint adopts on the joint law of value and sale.

\subsection{The three protocols}

Randomization is at the level of the pair and blocked by session, and both members of a pair receive the same protocol. Protocol F collects only final decisions, protocol S adds one intermediate measurement per trial, drawn between two possible moments and balanced over priors and intensities, and protocol R collects both intermediate measurements. A trial lasts 150 seconds, with a first request between seconds 10 and 40, the first certificate at second 50, a second request between seconds 60 and 90, the second certificate at second 105 and the final decision between seconds 120 and 150, so that the last intermediate window closes before the second certificate, a pause of thirty seconds that contains a new certificate separates the last measurement from the final window, and no window is extended for one protocol only. The length of the windows is fixed by the pilot, because the first hypothesis is tested on answers given within them and time pressure favors heuristics that ignore part of the available information, in the evidence surveyed by Spiliopoulos and Ortmann (2018), so that a tight deadline could produce the very neglect that the experiment measures. Following Henckel, Menzies, Moffatt and Zizzo (2022), who chose their time limit by piloting so that it would not bind subjects focused on the task, the pilot is run with windows of 45 seconds, and the experiment keeps windows of 30 seconds if 95 percent of the pilot answers to the intermediate requests arrive within that time, while otherwise it adopts windows of 45 seconds together with a shorter variant of sixteen purchase trials that drops the weak intensity, which carries about 6 percent of the information on the degree of neglect. The purchase task takes 55 minutes with the first calendar and 52 with the second, and a session lasts roughly 96 to 106 minutes by design.

\subsection{Controls and payment}

Eight lists with an announced value, in which the good is worth 36, 48, 66 or 78 points with certainty and the reservation price equals that value for both sellers, estimate on the same instrument the response noise, the attraction toward the center of the list and the individual tendency to understate, and sixteen choices between two urns with visible counts check the premise of the payment rule, that a participant prefers the higher probability of the same prize even when the difference is small. Each participant receives 8 euros for taking part and can win a single prize of 16 euros. One decision is drawn for payment, with probabilities of 48 percent for the final decisions, 13 percent for each of the two intermediate windows, 16 percent for the lists with an announced value and 10 percent for the urns, equal across protocols because an inactive slot keeps a reserve lottery at 50 percent, and the points earned in the drawn decision change the probability of the prize, which equals one half plus the points divided by 140, the steepest linear conversion that the two formats allow. A prize paid through a binary lottery keeps reservation prices free of risk attitudes if the participant reduces compound lotteries, an assumption that the design does not test and whose failure would distort the belief read from the first price far more than the degree of neglect, since a risk premium multiplies the odds by a constant, as the sale does, and the apparent degree of a buyer without neglect who values at 55 points a good worth 30 or 90 with equal chances stays below 0.01. The mechanism of Becker, DeGroot and Marschak punishes small departures from the true reservation price very lightly, a threshold five points above the optimum costing about 2.4 euro cents in expectation, and the design responds with the steep conversion, with unpaid practice lists after which the drawn price and the outcome are shown, and with the understatement term estimated on the lists with an announced value. Every programmed rule is declared and followed, so that the design contains no deception.

\section{Hypotheses and preregistered analysis}

\subsection{First hypothesis}

The first hypothesis is a positive mean degree of selection neglect in protocol S, estimated on the single intermediate measurement of each trial in the twelve joint trials and in the six single-seller trials together, so that the pull is estimated and removed. The estimand is the cross-person mean of the individual degree in the static restriction under the averaging rule, the test is one sided, and the primary procedure is an interval-censored likelihood built jointly on the thresholds, since a ratio of noisy prices is nonlinear and carries error in the variable used to build the benchmark. The declared measurement model writes the logit of the subjective posterior as a combination of the prior logit and of the log likelihood ratios of the certificates seen so far, with estimated weights and an individual intercept, lets the stated thresholds load on the center of the list through an attraction coefficient estimated from the randomization of the formats, subtracts an individual understatement, transports the error measured on the lists with an announced value into the purchase trials, and gives the individual intercept and the individual degree a bivariate normal distribution with free correlation. The primary test is the one-sided likelihood-ratio test of a zero mean degree at the level 0.025, with a critical value calibrated by parametric bootstrap under the null at the restricted estimates of the nuisance parameters. A studentized bootstrap over pairs, blocked by session, with 1,999 replications, is reported as a robustness check that does not rely on the parametric form of the measurement model, under two rules declared in advance, by which the bootstrap over pairs becomes primary if a goodness-of-fit check or the placebo check rejects the measurement model, and the calibrated test decides, with both results reported, if the checks pass and the two procedures disagree.

\subsection{Second hypothesis}

The second hypothesis is the equivalence between protocols F and R of the final price for the selective seller in the twelve joint trials, within a margin of three points, which is five percent of the distance between the two values. To avoid imputing a latent price, the analysis uses the observable policy index on the common economic range,
\[
P^{\mathrm{obs}} \;=\; 30+\sum_{p=31}^{90}\mathbf{1}\{\text{buy at price } p\},
\]
which uses exactly the same sixty decisions in the two formats and is defined for every policy. For each person the index minus the correctly conditioned benchmark at the objective final posterior is averaged over the twelve joint trials, the analysis compares the means of pairs with session fixed effects, and the two one-sided tests of Schuirmann (1987) each use level 0.025, so that equivalence requires the two-sided 95 percent interval to lie inside the interval from minus three to three and the family of the two primary claims is controlled conservatively at five percent. The margin is about a fifth of the gap between the benchmark prices averaged over the joint trials, 15.9 points, and it is wide for a small neglect, since a mean degree of 0.15 raises the same outcome by 2.4 points, a limit that is declared here and on which the judgment of an experimentalist is sought. An interval that crosses the margin is inconclusive about equivalence even when the difference from zero is not significant, and the result, where it obtains, is average equivalence of the protocols in this task and under these incentives, which does not prove that measuring is innocuous everywhere.

\subsection{Third hypothesis}

The third hypothesis concerns the form of the neglect and is examined only if the first is rejected. The null hypothesis is a rule weight of zero, that is, the averaging rule, against the alternative that the dampening rule has positive weight, tested by a likelihood ratio at 5 percent with a critical value calibrated by simulation on all the intensities. A rule fixed in advance decides whether the question is a preregistered hypothesis or an estimate reported with its confidence interval, the former holding only if the simulation shows a power of at least 80 percent to reject the averaging rule when the dampening rule is true and the mean degree is 0.25. The analysis with the rule weight is the one that the authors intend to implement if the time limits of the course project allow, and the analysis with the averaging rule alone is the fallback.

\subsection{Secondary analyses and diagnostics}

The secondary family, with a separate declared false discovery rate control, contains the independence, revelation and announced-value placebos, order effects, the comparison of protocol S with F and R, the persistence of coefficients between first and final measurement, the sensitivity to the center of the list and the urn choices. The main diagnostic concerns the attraction toward the center, since an attraction of ten percent produces, in a buyer who conditions correctly at the posterior 0.4 on a list centered at 55, apparent degrees of 0.102, 0.103 and 0.104 across the three intensities, nearly the common value that the neglect itself predicts, so that the restriction across intensities cannot separate the two accounts, and the lists with an announced value together with the random shift of the list are the instruments that do. Placebo trials predict a zero gap for every degree, revelation trials predict coinciding prices at the correctly updated value after availability is shown, and both predictions are testable without the structural model.

Two comparisons do not depend on the measurement model. For each pair, the mean gap between partners in the single-seller trials is set against the mean gap within persons in the joint trials, a positive difference on average indicating that stating the prices together compresses them. In the single-seller trials, moreover, the price that one partner states for the selective seller is set against the correct price at the belief revealed by the price that the other states for the always-available seller, after both prices have been corrected for the attraction measured by the two formats and for the understatement measured on the lists with an announced value, and a positive mean difference across pairs indicates neglect without the model of beliefs and of errors on which the likelihood rests. The correction is necessary, since in a preliminary simulation with an attraction of ten percent and no neglect the uncorrected comparison rejects at the 2.5 percent level in 97 percent of 1,000 samples of 50 pairs, against 3 percent with the correction, which detects a mean degree of 0.15 in 93 percent of the samples.

\section{Sample, pilot rule and feasibility}

The planning base is 300 persons in 150 pairs, fifty pairs per protocol, organized in ten main sessions of thirty, together with a separate pilot of 36 persons in two sessions of eighteen that does not enter the main analysis. Pairs are formed after the comprehension quiz among the participants admitted to the trials, and those who fail the quiz twice, together with one participant drawn at random when the number admitted is odd, receive the participation fee and do not take part. The base is conditional on an explicit rule declared before any data exist, since the central equivalence scenario reaches eighty percent power for the two one-sided tests at 300 only if the relevant standard deviation estimated in the pilot is about 5.92 or smaller, and the dossier states the freezing criterion that converts the pilot estimates into the final size.

The analysis procedure has been simulated in full for the design with two intensities and joint trials only, where at 100 subjects the asymptotic chi-square threshold rejects a true null in 3.2 percent of 2,000 replications, the critical value calibrated by simulation, 4.28, gives a level of 2.25 percent on independent null replications and a power of 99.5 percent at a degree of 0.15, the equivalence test reaches power 0.90 at a true difference of one point in the baseline scenario, and the pilot rule selects 300 in all 2,000 simulated pilots. Those numbers do not carry over to the present design without a new run, because that analysis would have attributed a positive pull to the neglect, and the simulation of the present design is in preparation under a plan fixed beforehand, with eight scenarios that combine a mean degree of 0, 0.15 and 0.25 with a pull of 0 and 0.10 and with the two rules of Section 4.3. The pilot also fixes the length of the windows and verifies comprehension, completion of the windows, censoring, sensitivity to the format and the difference in error between the lists with an announced value and the purchase trials, while the sign of any preliminary effect is never used to select the primary hypothesis. An eventual implementation would run as an online laboratory following the lab-like methodology developed by the CESARE group in Buso, De Caprariis, Di Cagno, Ferrari, Larocca, Lor\`e, Marazzi, Panaccione and Spadoni (2020), with subjects external to the course.

\section{What the design delivers and where it stops}

The design delivers a confirmatory core in prices with two preregistered hypotheses and a third whose status is decided by a declared rule, a measurement model that separates the neglect from the attraction toward the center of the list, from the pull between prices stated together and from an individual understatement, and trials that keep these separations testable without the model. It measures the neglect of a selection that is given to the buyer and does not respond to her price, and it is silent on the further inference required when the buyer's own offer decides whether the trade takes place and on behavior against a human seller. Every number in this document and in the dossier is a theoretical calculation, a benchmark diagnostic, a power scenario or a simulation of the analysis procedure, no human data have been collected, and the programs that travel with the design package reproduce the posterior enumeration, the benchmark prices under the two rules, the counterexample of the attraction toward the center, the equivalence scenarios, the preliminary check of the comparison between partners and the full simulation of the analysis for the design with two intensities, while the protocol, the explanatory note, this document and the programs are public at \url{https://github.com/francescopinna03/experimental_counterpart} and the replication package of the preprint at \url{https://github.com/francescopinna03/Dynamic_Identification}.

\begin{thebibliography}{99}

\bibitem{akerlof} G. A. Akerlof, The market for ``lemons'': quality uncertainty and the market mechanism, \emph{Quarterly Journal of Economics} 84 (1970), 488-500.

\bibitem{bdm} G. M. Becker, M. H. DeGroot, and J. Marschak, Measuring utility by a single-response sequential method, \emph{Behavioral Science} 9 (1964), 226-232.

\bibitem{brandts} J. Brandts and G. Charness, The strategy method: a survey of experimental evidence, \emph{Experimental Economics} 14 (2011), 375-398.

\bibitem{buso} I. M. Buso, S. De Caprariis, D. Di Cagno, L. Ferrari, V. Larocca, L. Lor\`e, F. Marazzi, L. Panaccione, and L. Spadoni, Lab-like findings of non-lab experiments: a methodological proposal and validation, CESARE Working Paper, LUISS Guido Carli, 2020.

\bibitem{casonplott} T. N. Cason and C. R. Plott, Misconceptions and game form recognition: challenges to theories of revealed preference and framing, \emph{Journal of Political Economy} 122 (2014), 1235-1270.

\bibitem{charnesslevin} G. Charness and D. Levin, The origin of the winner's curse: a laboratory study, \emph{American Economic Journal: Microeconomics} 1 (2009), 207-236.

\bibitem{danz} D. Danz, L. Vesterlund, and A. J. Wilson, Belief elicitation and behavioral incentive compatibility, \emph{American Economic Review} 112 (2022), 2851-2883.

\bibitem{dicagno} D. Di Cagno, W. G\"uth, T. Lohse, F. Marazzi, and L. Spadoni, Who cares when value (mis)reporting may be found out? An acquiring-a-company experiment with value messages and information leaks, \emph{Journal of Behavioral and Experimental Economics} 108 (2024), 102151.

\bibitem{enke} B. Enke, What you see is all there is, \emph{Quarterly Journal of Economics} 135 (2020), 1363-1398.

\bibitem{enkegraeber} B. Enke and T. Graeber, Cognitive uncertainty, \emph{Quarterly Journal of Economics} 138 (2023), 2021-2067.

\bibitem{espondavespa} I. Esponda and E. Vespa, Hypothetical thinking and information extraction in the laboratory, \emph{American Economic Journal: Microeconomics} 6 (2014), 180-202.

\bibitem{eysterrabin} E. Eyster and M. Rabin, Cursed equilibrium, \emph{Econometrica} 73 (2005), 1623-1672.

\bibitem{grether} D. M. Grether, Bayes rule as a descriptive model: the representativeness heuristic, \emph{Quarterly Journal of Economics} 95 (1980), 537-557.

\bibitem{harrison} G. W. Harrison, Theory and misbehavior of first-price auctions, \emph{American Economic Review} 79 (1989), 749-762.

\bibitem{henckel} T. Henckel, G. D. Menzies, P. G. Moffatt, and D. J. Zizzo, Belief adjustment: a double hurdle model and experimental evidence, \emph{Experimental Economics} 25 (2022), 26-67.

\bibitem{kagellevin} J. H. Kagel and D. Levin, The winner's curse and public information in common value auctions, \emph{American Economic Review} 76 (1986), 894-920.

\bibitem{pinna} F. Pinna, Dynamic identification of behavioral reasoning in finite games, SSRN Working Paper 7251718, version of August 8, 2026. \url{https://papers.ssrn.com/sol3/papers.cfm?abstract_id=7251718}.

\bibitem{samuelsonbazerman} W. F. Samuelson and M. H. Bazerman, The winner's curse in bilateral negotiations, in \emph{Research in Experimental Economics}, volume 3, JAI Press, 1985, 105-137.

\bibitem{schlag} K. H. Schlag, J. Tremewan, and J. J. van der Weele, A penny for your thoughts: a survey of methods for eliciting beliefs, \emph{Experimental Economics} 18 (2015), 457-490.

\bibitem{schuirmann} D. J. Schuirmann, A comparison of the two one-sided tests procedure and the power approach for assessing the equivalence of average bioavailability, \emph{Journal of Pharmacokinetics and Biopharmaceutics} 15 (1987), 657-680.

\bibitem{spiliopoulos} L. Spiliopoulos and A. Ortmann, The BCD of response time analysis in experimental economics, \emph{Experimental Economics} 21 (2018), 383-433.

\end{thebibliography}

\end{document}
